\section{Correctness and fair limits of the specified completion procedures}

\begin{annexitem}{annex:p12.fc.restart}{thm:p12-fc-restart}{Interreduction terminates and the fresh scan proves coverage}
\begin{proof}
Every current rule decreases the fixed reduction ordering. Thus
normalizing a finite term with the current finite rule set terminates;
no confluence assumption is needed to obtain a normal form.
During one interreduction phase there are finitely many rule records.
Temporarily give the two endpoints of each record persistent bookkeeping
identities, independent of which endpoint is currently written on the
left. Normalization replaces an endpoint only by a smaller term in
\(>\). Reorientation merely exchanges its displayed side. Deletion
removes a record, and no new record is introduced in this phase.
An infinite phase would therefore update some endpoint infinitely
often, giving an infinite descending sequence in \(>\). This is
impossible. At a processed nondeleted record at least one side was
reducible, so there is an actual decrease; a deleted record can be
deleted only once. The phase terminates or encounters an incomparable
normalized pair and reports failure.

Replacing a rule by its two normal forms under the other rules
preserves the generated equality: the retained rules equate each
old endpoint to its replacement, so the old and new equations imply
one another in the presence of those retained rules. If both normal
forms coincide, the old rule was already an equality consequence of
the others and may be deleted. A newly added critical equation is
an existing equational consequence of its peak. Thus every completed
interreduction and addition preserves the original theory.

For a fixed finite rule set, there are finitely many ordered rule
pairs and nonvariable positions. First-order unification supplies
the finite critical-pair list. If the final scan succeeds, its
normalization calculations exhibit an actual common reduct for
every pair. The critical-pair theorem supplies local confluence,
and the common reduction ordering supplies termination. Newman's
lemma gives confluence. The invariant and the normal-form theorem
therefore identify its equational theory with \(E_0\).
These arguments prove successful-output correctness, not that only
finitely many new rounds will be required.
\end{proof}
\end{annexitem}

\begin{annexitem}{annex:p12.fc.invariants}{thm:p12-fc-invariants}{All macrostep, phase and historical-rule invariants}
\begin{proof}
We first record consequences of the reduction-ordering hypothesis.
It is impossible that \(a>b\) and \(b=C[a\sigma]\): repeated
substitution and context closure would give the infinite descending
sequence
\[
 a>C[a\sigma]>C[C\sigma[a\sigma^2]]>\cdots.
\]
In particular a term cannot exceed a context containing itself.
An oriented pair \(\ell>r\) has no variable in \(r\) absent
from \(\ell\). If such a variable \(x\) existed, substituting
\(\ell\) for \(x\) and fixing variables of \(\ell\) would
give either \(\ell>\ell\) or \(\ell>C[\ell]\), both
impossible. Nor can \(\ell\) be a variable: if it occurs in
\(r\), equality or a proper self-context contradicts the order;
if it does not occur, substituting \(r\) for it contradicts
irreflexivity. Thus the inserted orientation is an ordinary valid rule.

Induct on iterations to show all current rules decrease \(>\).
The new rule does so by its selection. Old retained right sides
are reduced by \(S_i=R_i\cup\{\ell\to r\}\), whose rules
all decrease \(>\), so their old left sides still exceed their
new right sides by transitivity. Trivial deletion and marking do
not change rules. It follows also that every rule ever generated
decreases the same ordering; both \(R_*\) and \(R_\infty\)
are terminating.

We next prove interreduction at each stage. It holds initially.
At insertion, \(\ell\) and \(r\) are \(R_i\)-normal by
the equation-normalization step. The term \(r\) cannot contain
an instance of \(\ell\), since \(\ell>r\) and the first
paragraph prohibit that. Hence \(r\) is also normal under the
new rule. Every retained old left side is irreducible by the new
rule by the definition of \(K\), and by every other old rule
by induction. The new left side is irreducible by all old rules.
Changing right sides does not change whether a left side matches
a subterm. Finally, the retained right-side normal forms are
\(S_i\)-normal and the resulting left sides are a subset of
the old left sides together with \(\ell\). They are therefore
normal under the resulting rules. This proves both required
interreduction conditions. The other iterations preserve them.

For termination of an equation phase, take its finite collection
of pending-equation and rule records. As bookkeeping only, attach
two endpoint slots to each record. A selected pending record is
deleted, or becomes the new rule record after its endpoints have
been normalized. A collapsed old rule becomes a pending record
with one endpoint strictly reduced. Retained rules keep their
records and normalize one endpoint. Reorientation exchanges which
slot is displayed on which side. Duplicate equations can delete
records but introduce none. Thus the phase has at most the original
finite number of slots, and every changed slot decreases in \(>\).
There can be only finitely many actual slot decreases, since otherwise
one of the finitely many slots would decrease infinitely often.
After the last decrease there can be no further collapse. Every
remaining equation iteration therefore deletes a trivial equation
or moves an unchanged pending pair into the rule set. Both decrease
the number of pending records. The phase ends in finitely many
iterations in a nonfailure run, independently of equation selection.

An equation present in a phase remains pending until selected;
only the selected equation is removed by a macrostep. Since the
phase ends empty, each such equation is selected. Its two sides
reduce in \(R_i\) to equal normal forms, or one of those normal
forms reduces to the other by the newly inserted rule. These are
finite \(R_*\)-reductions to a common term. This proves the
claimed joining property for every pending equation, including
critical equations and equations created by collapse.

For soundness, induct on generated data. Initial equations are
consequences of \(E_0\). A selected equation is normalized by
already justified rules, so its new orientation is still an
\(E_0\)-equality. A collapsed equation equates a reduct of an
old left side with its old right side; the old rule and the new
rule justify both equalities. A retained right-side normalization
uses old rules and the justified new rule. Critical pairs are
the two reducts of one term, so are also justified equations.
Thus every rule of \(R_*\) is an \(E_0\)-consequence.
Conversely every equation of \(E_0\) has an \(R_*\)-join by
the preceding paragraph. Taking congruence closure gives
\(\equiv_{R_*}=\equiv_{E_0}\).

Finally, a label is created once and never reused. Until deletion
its left side is fixed, and its successive proper right-side changes
are finite \(R_*\)-reductions and therefore strict decreases in
\(>\). There are only finitely many such changes. If the label
is never deleted, its final right side is present unchanged from
some stage onward. Every earlier right-side version reduces to
that final version in finitely many \(R_*\)-steps. The same
statement holds up to the deletion stage for a deleted label.
\end{proof}
\end{annexitem}

\begin{annexitem}{annex:p12.fc.containment}{thm:p12-fc-containment}{A concrete rank for strict pattern containment}
\begin{proof}
Containment is reflexive and transitive: compose the substitutions
and the contexts witnessing the two containments. If \(a\)
contains \(b\), then \(\lvert a\rvert\geq\lvert b\rvert\),
because substituting a finite term for a variable cannot decrease
node count and passing to a superterm cannot decrease it.
If these sizes are equal, the occurrence must be at the root, and
every variable occurring in \(b\) is replaced by a single node,
which may be a variable or a constant.
Then \(\#\operatorname{Var}(a)\leq\#\operatorname{Var}(b)\).
If the variable counts also agree, no occurring variable can have
been replaced by a constant, and that finite variable map is
injective on the occurring variables and can be inverted there;
the terms are variable-renaming variants and contain one another.
Consequently \(a\gg b\) strictly decreases the lexicographic rank
\[
 \bigl(\lvert a\rvert,
       \lvert a\rvert-\#\operatorname{Var}(a)\bigr)
 \in\N^2.
\]
Both coordinates are nonnegative. The first can decrease only
finitely often in a descending rank sequence; once it stabilizes,
the second cannot decrease infinitely. Thus \(\gg\) is well-founded.

If a new rule with left side \(\ell\) deletes the old left
side \(a\), then \(a\) contains \(\ell\) by the definition
of \(K\). But \(\ell\) is normal under the old rule set,
so it cannot contain \(a\), a current left-hand side. Hence
\(a\gg\ell\), including root-instance collapses and proper
subterm collapses.

At the next iteration, a retained label still has the same left
side. A deleted label is replaced, for applicability purposes, by
the new left side contained in the old one. Any term containing
the old pattern contains the new pattern by transitivity, and is
therefore still reducible. Trivial equation deletion and marking
do not affect reducibility. Induction over later iterations proves
persistence of reducibility.
\end{proof}
\end{annexitem}

\begin{annexitem}{annex:p12.fc.fair-coverage}{thm:p12-fc-fair-coverage}{Coverage of persistent overlaps through historical versions}
\begin{proof}
Consider the two persistent rule applications at \(a\).
If their occurrences are disjoint or one lies inside a substituted
variable of the other, the corresponding cases of
Annex~\ref{annex:p12.rewriting.critical-pairs} join the endpoints
using only those persistent rules. This includes nonlinear copying
and erasure. Take \(d_i=b_i\) and that common reduct for \(e\).
The required first steps belong to \(R_*\), since
\(R_\infty\subseteq R_*\).

For a nonvariable overlap, let the persistent labels be \(j\leq k\).
Label \(j\) was created no later than \(k\), and neither is
ever deleted. Fairness therefore selects \(k\) at some
critical-pair iteration when \(j\) is also present. Their
left sides are already the eventual left sides. Let their right
sides at that iteration be \(r'_1,r'_2\), and their final
right sides be \(r_1,r_2\). Each \(r'_i\) reduces to \(r_i\)
in \(R_*\) by the preceding invariant lemma.

Apply the historical versions at the same occurrences and with the
same matches used by the persistent peak. This gives
\(a\to_{R_*}d_i\) and \(d_i\to_{R_*}^*b_i\).
The two historical reducts are the contexted substitution instance
of the critical pair computed for those left sides. That pair
was put into the pending set, in the appropriate overlap direction;
self-overlaps are included when \(j=k\). It has a common
\(R_*\)-reduct by the pending-equation invariant. Context and
substitution closure lift the two joining reductions to a common
term \(e\) for \(d_1,d_2\). This proves the displayed diagram
without presuming a join of the final right-side versions.
\end{proof}
\end{annexitem}

\begin{annexitem}{annex:p12.fc.limit}{thm:p12-fc-limit}{The complete persistent-rule transfer induction}
\begin{proof}
Write \(\to_*\) for one-step rewriting by \(R_*\), and
\(\to_\infty\) for one-step rewriting by \(R_\infty\), only
within this proof. Their reflexive-transitive closures carry a
superscript \(*\). Every one-step reduction of either kind
strictly decreases the fixed ordering \(>\).

For each term \(a\), prove the following three assertions, in the
listed order, by well-founded induction on \(a\):
\begin{enumerate}
\item[(A)] If \(a\to_*^*b\), some \(p\) satisfies
      \(a\to_\infty^*p\) and \(b\to_\infty^*p\).
\item[(B)] If \(a\to_\infty^*b_1,b_2\), then
      \(b_1,b_2\) have a common \(R_\infty\)-reduct.
\item[(C)] If \(a\to_*^*b_1,b_2\), then
      \(b_1,b_2\) have a common \(R_\infty\)-reduct.
\end{enumerate}
Assume all three assertions for every strictly smaller term.

\paragraph{Assertion (A).}
For a zero-step path take \(p=a=b\). For a positive path write
\(a\to_*a_1\to_*^*b\). Its first step uses a historical rule
with label \(k\), left side \(\ell\), and right side \(r\).
For this fixed \(a\), use a subsidiary well-founded induction on
\(\ell\) under \(\gg\), asserting the result for every path
whose first rule has that left-side rank.

If label \(k\) is never deleted, it has a final version
\(\ell\to\bar r\) in \(R_\infty\), with
\(r\to_*^*\bar r\). At the original occurrence and substitution
this gives \(a\to_\infty a_2\) and
\(a_1\to_*^*a_2\). Since \(a>a_1\), assertion (C) at
\(a_1\), applied to its reducts \(a_2,b\), supplies a common
persistent reduct \(p\). Then
\(a\to_\infty a_2\to_\infty^*p\) and
\(b\to_\infty^*p\), proving (A) in this case.

If the label is eventually deleted, its version just before deletion
is \(\ell\to r'\), with \(r\to_*^*r'\). The new rule
\(h\to q\) performs a step \(\ell\to\ell'\), and the
equation \(\ell'\approx r'\) is made pending. Here
\(\ell\gg h\) by Lemma~\ref{thm:p12-fc-containment}.
Lift this step to the original context and substitution, obtaining
\(a\to_{\{h\to q\}}a'_1\). The changed old right side
gives \(a_1\to_*^*a_2\). The pending equation has an
\(R_*\)-join, which lifts to reductions
\(a'_1\to_*^*d\) and \(a_2\to_*^*d\).
Thus \(a_1\) reduces in \(R_*\) to both \(b\) and \(d\).
Assertion (C) at the smaller term \(a_1\) gives
\(b\to_\infty^*b'\) and \(d\to_\infty^*b'\).
Now the path
\[
 a\to_{\{h\to q\}}a'_1\to_*^*d\to_\infty^*b'
\]
has a first rule with strictly smaller containment rank. The
subsidiary induction gives \(a\to_\infty^*p\) and
\(b'\to_\infty^*p\) for some \(p\). Concatenating from
\(b\) proves (A). Both label-lifetime cases are now covered.

\paragraph{Assertion (B).}
If either persistent path has length zero, the other endpoint is
a common persistent reduct. Otherwise write the first steps as
\(a\to_\infty a_1\to_\infty^*b_1\) and
\(a\to_\infty a_2\to_\infty^*b_2\).
The fair-coverage lemma supplies historical one-step reducts
\(d_1,d_2\) of \(a\) and a common term \(e\), with
\(d_i\to_*^*a_i\) and \(d_i\to_*^*e\).
Both \(d_i\) are strictly below \(a\). Assertion (C) at
\(d_1\), applied to \(b_1,e\), gives
\(b_1\to_\infty^*u\) and \(e\to_\infty^*u\).
Since \(d_2\to_*^*e\to_\infty^*u\), assertion (C) at
\(d_2\), applied to \(b_2,u\), gives
\(b_2\to_\infty^*p\) and \(u\to_\infty^*p\).
This joins \(b_1,b_2\), proving (B).

\paragraph{Assertion (C).}
Apply the already proved (A) at \(a\) separately to its two
\(R_*\)-reducts. Obtain persistent paths from \(a\) to
\(p_1,p_2\), and from \(b_i\) to \(p_i\).
Assertion (B) at \(a\) joins \(p_1,p_2\), and concatenation
joins \(b_1,b_2\). This proves (C). There is no circular use:
(A) and (B) used (C) only at strictly smaller terms; (C) uses
the already established (A) and (B) at the current term.

The induction proves all three assertions. In particular (B)
gives confluence of \(R_\infty\), while (C) also gives
confluence of \(R_*\). Both are terminating by the common order.
For every rule of \(R_*\), (A) makes its two sides equal
modulo \(R_\infty\). Hence
\(\equiv_{R_*}\subseteq\equiv_{R_\infty}\); the reverse
inclusion follows from \(R_\infty\subseteq R_*\).
The generated-equation invariant identifies both with
\(\equiv_{E_0}\).

Finally two persistent rules and their final right sides coexist
unchanged at every sufficiently late iteration. A persistent
right-side redex, or a left-side redex created by a different
persistent label, would therefore contradict interreduction at
that finite iteration. Thus the limit is interreduced as well.
\end{proof}
\end{annexitem}

\begin{annexitem}{annex:p12.fc.computation}{thm:p12-fc-computation}{Successful outputs, finite limits and equality semidecision}
\begin{proof}
A successful finite run, extended constantly, satisfies fairness:
each surviving label is marked, hence was selected, and each other
created label was deleted. Its limit is the final finite rule set,
so the preceding theorem proves it canonical for \(E_0\).

We first derive one useful consequence of the transfer assertion.
If \(a\to_{R_*}b\), then \(a\) is reducible by
\(R_\infty\). Otherwise (A) forces its common persistent
reduct with \(b\) to be \(a\) itself. This would give
\(b\to_{R_\infty}^*a\), contradicting the strict decrease
\(a>b\) and the fact that every persistent step decreases
the same ordering. Thus reducibility by the union and by the
persistent limit are equivalent.

Suppose \(R_\infty\) is finite. Choose a stage after every one
of its final labeled records is present forever. At this and all
later stages, any extra current rule would have a left side
reducible in \(R_*\), hence by \(R_\infty\). A different
persistent rule would then reduce that current left side,
contradicting interreduction. If it were the same label, uniqueness
of a label's current record makes it the persistent record, not
an extra rule. Therefore the current rule set equals
\(R_\infty\) at every such stage.
No further insertion can occur. The finite equation phase ends,
and the finitely many unmarked labels are marked one by one,
with only finite equation phases in between. The procedure stops
with success. Its contrapositive gives the infinite-limit assertion.

For a query pair, define \(a_i,b_i\) by normalizing the preceding
query values with the current finite \(R_i\). All resulting
reductions use \(E_0\)-equalities, so an observed equality
\(a_i=b_i\) is sound. If the run is infinite, each query term
can change only finitely many times, by well-foundedness of \(>\).
Let its eventual values be \(a',b'\). Neither has an
\(R_\infty\)-redex: the rule supplying such a redex would be
present at every sufficiently late iteration, contradicting
normalization of an unchanging query value. They are therefore
the unique persistent normal forms of the original terms.
If the terms are \(E_0\)-equal, these normal forms coincide,
and equality is detected at a finite iteration. If they are not
equal, soundness prevents acceptance. At a successful finite halt,
the final finite rule set computes both persistent normal forms
and comparison also decides inequality. None of these arguments
applies after the incomparable-equation failure exit.
\end{proof}
\end{annexitem}

\begin{annexitem}{annex:p12.fc.antichain}{thm:p12-fc-antichain}{Persistent left sides and the infinite-antichain consequence}
\begin{proof}
The persistent limit is interreduced. Take two distinct persistent
rule records with left sides \(a,b\). If \(a\) contains \(b\),
the rule with left side \(b\) reduces \(a\) at the witnessing
occurrence and substitution. This contradicts irreducibility of
\(a\) under the other persistent rules. Interchanging \(a,b\)
excludes the reverse containment. In particular their left sides
cannot be equal or variable-renaming variants; either case would
give mutual containment. There is therefore at most one persistent
rule for each left-side pattern, and the left sides of an infinite
limit form an infinite containment antichain.

The finite-limit theorem shows that an infinite nonfailure fair run
has an infinite limit. Its left-side antichain is therefore infinite.
If the term language admits no such antichain, an infinite run under
these hypotheses is impossible. The only remaining nonfailure exit
is success, since all equation phases and every individual finite
critical-pair computation terminate. Neither this argument nor its
pattern-containment relation includes matching modulo a further
equational congruence; the criterion has only the stated ordinary
term-language scope.
\end{proof}
\end{annexitem}

\begin{annexitem}{annex:p12.fc.scheduler}{thm:p12-fc-scheduler}{Fairness of label age and the explicit finite run}
\begin{proof}
Fix an allocated label \(k\). If it is deleted, its fairness
requirement is satisfied. Otherwise, while it remains unmarked,
the least-label scheduler can choose before it only unmarked
labels smaller than \(k\). There are finitely many of those;
new labels are greater, and a selected old label remains marked
unless deleted. Each intervening equation phase terminates.
Consequently \(k\) is selected after finitely many critical-pair
iterations. This proves fairness without any assumption that
syntactic size alone has finitely many competitors.

Use the ordering interpreting \(a=2,c=1,b=0,f(n)=n+1\).
On this unary signature each term evaluates either to a constant
integer or to one variable plus an integer. Universal comparison
over natural-number assignments is therefore decidable, so the
example meets the algorithm's recursive-order hypothesis.
Select \(f(a)\approx b\) first. Its sides are already normal
and ordered, giving unmarked label 1, \(f(a)\to b\).
Select \(a\approx c\) next, giving label 2, \(a\to c\).
It reduces the old left side \(f(a)\) to \(f(c)\), so label 1
is deleted and \(f(c)\approx b\) becomes pending. Its two
sides are normal under label 2 and ordered, so insert label 3,
\(f(c)\to b\).

The pending set is now empty. Select label 2: its self-overlap
has identical sides \(c,c\). Mark it and discard that trivial
equation. Select label 3: its self-overlap gives \(b,b\), and
there is no overlap with label 2, since the only nonroot symbol
in \(f(c)\) is \(c\), not \(a\). Both root directions
between the distinct left sides also fail to unify. Mark label 3
and discard its trivial self-equation. All records are marked
and the pending set is empty, so success is reported.
The complete scan, not an assumed final joinability condition,
has discharged every pair required for this returned rule set.
\end{proof}
\end{annexitem}

\begin{annexitem}{annex:p12.fc.associativity}{thm:p12-fc-associativity}{Ordering, every overlap and canonical leaf sequences}
\begin{proof}
Write the binary operation in infix notation. Fix any strict total
precedence on this finite signature. Give its binary symbol,
its constants and the variable weight the positive weight one.
These data satisfy the original Knuth--Bendix admissibility
conditions of Definition~\ref{def:p12-ro-kbo-weights}; no zero-weight
unary symbol occurs. Both sides of the associativity rule have
weight five and the same count of each variable. Their outer
symbols agree. In their first arguments, \(x\cdot y\) has
weight three and exceeds \(x\), of weight one; its count of
each variable is at least the corresponding count in \(x\).
The historical left-to-right equal-weight branch therefore orders
the whole left side above the right side. Theorem~\ref{thm:p12-ro-kbo-termination}
makes this a reduction ordering, proving termination.

The left side has just two nonvariable positions: the root and
the left child. A root overlap of its renamed copies gives an
identical pair. The left-child overlap has peak term
\(((x\cdot y)\cdot z)\cdot w\). Its root-first branch is
\[
 ((x\cdot y)\cdot z)\cdot w
 \longrightarrow (x\cdot y)\cdot(z\cdot w)
 \longrightarrow x\cdot(y\cdot(z\cdot w)).
\]
Its inner-first branch is
\[
 ((x\cdot y)\cdot z)\cdot w
 \longrightarrow (x\cdot(y\cdot z))\cdot w
 \longrightarrow x\cdot((y\cdot z)\cdot w)
 \longrightarrow x\cdot(y\cdot(z\cdot w)).
\]
There are no other rule pairs or nonvariable overlap positions.
The critical-pair criterion and Newman's lemma give confluence.
The rule is precisely the initial associativity equation, so its
equational theory is that initial theory.

Its right side is irreducible, and its left side has no proper
redex under the single rule. The restart procedure leaves it
unchanged at interreduction and accepts the joining critical-pair
scan. Huet's procedure inserts its one unmarked label, computes
the identical root pair and the displayed joining pair, deletes
their coincident normalized equations, and stops with that label
marked. Thus neither procedure inserts a further rule.

Every reduction preserves the sequence of leaves read from left
to right: the local replacement moves only parentheses. Conversely
a normal term for this rule has no binary-rooted left child.
Induction on its finite tree shows it is exactly the fully
right-associated product of its leaf sequence, or the one leaf
itself. Hence each leaf sequence determines exactly one normal
form. Distinct sequences have distinct normal forms and cannot be
identified by the associativity theory; identical sequences normalize
to the same term.
\end{proof}
\end{annexitem}

\begin{annexitem}{annex:p12.fc.permutation-failure}{thm:p12-fc-permutation-failure}{The variable-swap contradiction and the meaning of failure}
\begin{proof}
Suppose \(f(x,y)>f(y,x)\). Apply the substitution exchanging
\(x\) and \(y\) and fixing every other variable. Substitution
stability gives \(f(y,x)>f(x,y)\). Transitivity would yield
\(f(x,y)>f(x,y)\), contradicting strictness. Starting with the
opposite comparison produces the same contradiction. The two
terms are syntactically distinct, so they are neither comparable
nor equal.

Initially there are no rules, so the two terms are their own
current normal forms. Both specified procedures therefore take
their incomparable-equation failure exit. The equation remains an
input axiom throughout; this exit proves no inequality in its
equational theory. For example, in the language consisting only of
this binary symbol and constants, the corresponding word problem
is decidable by a direct bottom-up canonicalization: retain each
leaf, canonicalize the two children of each binary node, and put
their canonical encodings in a fixed decidable lexicographic order.
The resulting encoding is invariant under a swap at any node.
Conversely, equal encodings imply by induction that the children
can be paired, after at most one swap, with equal encodings;
recursively swap inside each child to get an equational conversion.
Thus equal encodings are exactly equality modulo commutativity.
This is a separate decision procedure; it does not orient the
universal commutativity axiom by a substitution-stable ordering.
\end{proof}
\end{annexitem}
